%% bk-howto.tex  前付「この本の使い方」（2ページ）。第2巻から共通。  \bkHowToPages
%% 第3回編集委員会：子どもに要る情報だけにする。★の表（巻ごとの中身の一覧）と路線マークの一覧（4・5ページと重なる）はのせない。
%% 問題はノートに解くので、ノートの使い方を足した。
\makeatletter
\newcommand\bk@howstep[7]{% x, y, アイコン, 色, 名前, 説明, ノード名
  \node[draw=#4, line width=1pt, fill=white, rounded corners=2.5mm,
    minimum width=30.5mm, minimum height=40mm, anchor=north] (#7) at (#1,#2) {};% 第11回：2行の名前にルビがあっても説明が箱からはみ出さない高さ（第7回の37 mm から）
  \node[anchor=north] (#7i) at ([yshift=-2.5mm]#7.north) {\bkicon[1.7em]{#3}{#4}};
  \node[anchor=north, font=\sffamily\bfseries\small, text=#4, inner sep=0pt, align=center,
    minimum height=12.5mm, text width=29.5mm] (#7l) at ([yshift=-1mm]#7i.south) {#5};% 第11回：ルビつき2行の名前が入る高さ
  \node[anchor=north, font=\sffamily\footnotesize, text width=28mm, align=center, text=bkInk, inner sep=0pt]
    at ([yshift=-2.0mm]#7l.south) {#6};}

\newcommand\bkHowToPages{%
\clearpage\setcounter{page}{6}%
\global\bk@decofalse\bkfolio{true}\bksetrunning{}{}{}%
\noindent{\sffamily\bfseries\fontsize{24}{30}\selectfont\color{bkDeep}この\R{本}{ほん}の\R{使}{つか}い\R{方}{かた}}\par\vspace{1mm}
{\color{bkVol}\hrule height 1.4pt}\vspace{4mm}

この\R{本}{ほん}は、\R{数学}{すう|がく}を「\R{物理}{ぶつ|り}や\R{工学}{こう|がく}を\R{考}{かんが}えるための\R{道具}{どう|ぐ}」として\R{学}{まな}ぶ\R{本}{ほん}です。ボールが\R{落}{お}ちる、お\R{湯}{ゆ}が\R{冷}{さ}める、\R{振}{ふ}り\R{子}{こ}がゆれる。そんな\R{身近}{み|ぢか}なできごとを、\R{数}{かず}と\R{式}{しき}で\R{書}{か}きあらわしてみませんか？

\bksubsection{1}{1つの\R{章}{しょう}は、いつも\R{同}{おな}じ\R{順番}{じゅん|ばん}}

\begin{bkfigure}
\begin{tikzpicture}[x=1mm, y=1mm]
  \bk@howstep{0}{0}{hatena}{bkYellow!85!black}{はてな？}{\R{身近}{み|ぢか}なふしぎから\\\R{始}{はじ}めよう}{s1}
  \bk@howstep{35}{0}{yattemiyou}{bkGreen}{やってみよう}{\R{手}{て}を\R{動}{うご}かして\\\R{調}{しら}べよう}{s2}
  \bk@howstep{70}{0}{doguu}{bkNavy}{\R{数学}{すう|がく}の\R{道具}{どう|ぐ}}{\R{新}{あたら}しい\R{考}{かんが}え\R{方}{かた}を\\\R{学}{まな}ぼう}{s3}
  \bk@howstep{105}{0}{reidai}{bkBlue}{\R{例題}{れい|だい}}{\R{解}{と}き\R{方}{かた}のお\R{手本}{て|ほん}を\\\R{見}{み}よう}{s4}
  \bk@howstep{0}{-48}{jibun}{bkOrange}{じぶんでとこう}{お\R{手本}{て|ほん}をまねして、\\ノートに\R{解}{と}こう}{s5}
  \bk@howstep{35}{-48}{tsumazuki}{bkRed}{まちがえやすい\\ところ}{まちがいの\R{理由}{り|ゆう}を\\\R{知}{し}ろう}{s6}
  \bk@howstep{70}{-48}{butsuri}{bkPurple}{\R{物理}{ぶつ|り}・\R{工学}{こう|がく}\\ではこう\R{使}{つか}う}{\R{何}{なに}の\R{役}{やく}に\R{立}{た}つか\\\R{知}{し}ろう}{s7}
  \bk@howstep{105}{-48}{check}{\bkTF}{\R{章末問題}{しょう|まつ|もん|だい}}{\R{章}{しょう}のまとめに\\チャレンジ}{s8}
  \foreach \a/\b in {s1/s2,s2/s3,s3/s4,s5/s6,s6/s7,s7/s8}
    {\draw[bk arrow, bkGray] (\a.east) -- (\b.west);}
  \draw[bk arrow, bkGray, rounded corners=2mm, shorten <=1.5mm, shorten >=1.5mm] (s4.south) -- ++(0,-4) -| (s5.north);% 第7回：2段のあいだを広げ、矢印を箱から1.5 mm はなす
\end{tikzpicture}
\end{bkfigure}
\vspace{2.5mm}

\R{例題}{れい|だい}と「じぶんでとこう」で、\R{章末問題}{しょう|まつ|もん|だい}に\R{出}{で}てくる\R{問題}{もん|だい}の\R{形}{かたち}を、ぜんぶ1\R{回}{かい}ずつ\R{見}{み}せています。こまったら、\R{例題}{れい|だい}にもどろう。

\bksubsection{2}{\R{章}{しょう}のはじめの「\R{道}{みち}しるべ」}

\leavevmode\R{章}{しょう}の\R{見出}{み|だ}しのすぐ\R{下}{した}に、こんな\R{行}{ぎょう}があります。
\par\vspace{2mm}
\noindent\begin{minipage}{\linewidth}\bkmichishirube{1-14,2-2}{3-5,5-9}\end{minipage}
\R{前}{まえ}に\R{読}{よ}んでおくとよい\R{章}{しょう}と、この\R{章}{しょう}の\R{考}{かんが}えが\R{育}{そだ}っていく\R{章}{しょう}です。わからない\R{言葉}{こと|ば}や\R{式}{しき}が\R{出}{で}てきたら、\R{前}{まえ}の\R{章}{しょう}にもどって\R{確}{たし}かめよう。

\clearpage
\bksubsection{3}{★は「どこまで\R{登}{のぼ}るか」のしるし}

\noindent\begin{minipage}[c]{0.44\linewidth}
\begin{tikzpicture}[x=1mm, y=1mm]
  % 山を5つの高さに分ける（左に★の数）
  \foreach \k/\c in {1/14,2/28,3/44,4/62,5/82} {
    \pgfmathsetmacro\bk@yb{(\k-1)*7.4}\pgfmathsetmacro\bk@yt{\k*7.4}
    \fill[bkTheme!\c] ({38-22*(1-\bk@yb/37)},\bk@yb) -- ({38+22*(1-\bk@yb/37)},\bk@yb)
      -- ({38+22*(1-\bk@yt/37)},\bk@yt) -- ({38-22*(1-\bk@yt/37)},\bk@yt) -- cycle;
    \node[anchor=east, inner sep=0pt] at (13,{(\k-0.5)*7.4}) {\bkstarmeter{\k}{\k}};}
  \draw[bkInk!70, line width=0.8pt] (38,37) -- (38,44);
  \fill[bkRed] (38,44) -- (44,42) -- (38,40) -- cycle;
\end{tikzpicture}
\end{minipage}\hfill
\begin{minipage}[c]{0.52\linewidth}
\R{章}{しょう}と\R{小見出}{こ|み|だ}しには★がついています。★の\R{数}{かず}は、\R{内容}{ない|よう}の\R{高}{たか}さ。\R{山登}{やま|のぼ}りにたとえると、★1つは「はじめの\R{一歩}{いっ|ぽ}」、★5つは「\R{頂上}{ちょう|じょう}」です。

\R{白}{しろ}ぬきの☆は「この\R{章}{しょう}の\R{中}{なか}で、ここまで\R{登}{のぼ}る」というしるし。\R{自分}{じ|ぶん}の\R{読}{よ}める★まで\R{読}{よ}んで、\R{先}{さき}へ\R{進}{すす}んでもだいじょうぶ。
\end{minipage}

\bksubsection{4}{ノートを1\R{冊}{さつ}\R{用意}{よう|い}しよう}

\noindent\begin{minipage}[c]{0.5\linewidth}
「じぶんでとこう」と\R{章末問題}{しょう|まつ|もん|だい}は、ノートに\R{解}{と}きます。\R{本}{ほん}には\R{書}{か}きこまないので、\R{何度}{なん|ど}でもやり\R{直}{なお}せます。

\R{問題}{もん|だい}の\R{番号}{ばん|ごう}を\R{書}{か}き、\R{図}{ず}をかいてから\R{式}{しき}を\R{書}{か}こう。\R{答}{こた}えには\R{単位}{たん|い}をつけます。
\end{minipage}\hfill
\begin{minipage}[c]{0.46\linewidth}
\begin{tikzpicture}[x=1mm, y=1mm]
  \fill[bkLine!60] (1.2,-1.2) rectangle (61.2,44.8);
  \draw[bkGray, line width=0.6pt, fill=white] (0,0) rectangle (60,46);
  \foreach \y in {7,14,...,42} \draw[bkBlue!25, line width=0.4pt] (2,\y) -- (58,\y);
  \draw[bkRed!40, line width=0.5pt] (10,0) -- (10,46);
  \node[anchor=base west, font=\sffamily\bfseries\footnotesize, text=\bkTT] at (1.5,36.2) {\R{例}{れい}};
  \node[anchor=base west, font=\footnotesize, text=bkInk] at (12,36.2) {じぶんでとこう 1};
  % 図
  \draw[bkInk!70, line width=0.6pt] (13,25) -- (49,25);
  \foreach \x in {13,19,...,49} \draw[bkInk!70, line width=0.5pt] (\x,24) -- (\x,26);
  \draw[bk arrow, bkTheme] (13,28.5) -- (49,28.5);
  \node[anchor=base west, font=\footnotesize, text=bkInk] at (12,15.2) {$\qty{6}{m/s}\times\qty{15}{s}=\qty{90}{m}$};
  \node[anchor=base west, font=\footnotesize\bfseries, text=bkInk] at (12,8.2) {\R{答}{こた}え\ \qty{90}{m}};
\end{tikzpicture}
\end{minipage}

\bksubsection{5}{この\R{本}{ほん}の3つのやくそく}

\begin{enumerate}[label={\bkcircled[bkNavy]{\arabic*}}, leftmargin=2.4\zw, itemsep=3pt]
\item \R{数}{かず}には、いつも\R{単位}{たん|い}をつけよう。「5」ではなく「\qty{5}{m}」と\R{書}{か}きます。
\item \R{式}{しき}を\R{書}{か}く\R{前}{まえ}に、\R{表}{ひょう}や\R{図}{ず}をかいて\R{考}{かんが}えよう。
\item くり\R{返}{かえ}しの\R{計算}{けい|さん}は、\R{電卓}{でん|たく}やコンピュータにまかせてもだいじょうぶ。
\end{enumerate}

\R{問題}{もん|だい}の\R{答}{こた}えは、\R{本}{ほん}の\R{後}{うし}ろの「こたえ」にあります。
}

%% ------------------------------------------------------------
%% 第1巻の「この本の つかいかた」（p.8〜9。第5回編集委員会）
%%   1年生が ひとりで よめるように、ひらがな・分かち書き・大きな字。書きこむ本であることを見せる。
%% ------------------------------------------------------------
\newcommand\bk@howstepone[7]{% x, y, アイコン, 色, 名前, 説明, ノード名
  \node[draw=#4, line width=1pt, fill=white, rounded corners=2.5mm,
    minimum width=31mm, minimum height=45mm, anchor=north] (#7) at (#1,#2) {};
  \node[anchor=north] (#7i) at ([yshift=-2.5mm]#7.north) {\bkicon[1.8em]{#3}{#4}};
  \node[anchor=north, font=\sffamily\bfseries\fontsize{11.5}{14}\selectfont, text=#4, inner sep=0pt, align=center,
    minimum height=10mm, text width=30mm] (#7l) at ([yshift=-1mm]#7i.south) {#5};
  \node[anchor=base, font=\fontsize{10.5}{16}\selectfont, text width=28mm, align=center, text=bkInk, inner sep=0pt]
    at ([yshift=-4.4mm]#7l.south) {#6};}% 第5回：ルビのある説明も1行目の高さをそろえる
\newcommand\bkHowToPagesOne{%
\clearpage\setcounter{page}{8}%
\begingroup\bkvolonebody
\global\bk@decofalse\bkfolio{true}\bksetrunning{}{}{}%
\noindent{\sffamily\bfseries\fontsize{24}{30}\selectfont\color{bkDeep}この \R{本}{ほん}の つかいかた}\par\vspace{1mm}
{\color{bkVol}\hrule height 1.4pt}\vspace{3mm}

\bksubsection{1}{1つの \R{章}{しょう}は、 いつも おなじ じゅんばん}

\begin{bkfigure}
\begin{tikzpicture}[x=1mm, y=1mm]
  \bk@howstepone{0}{0}{hatena}{bkYellow!85!black}{はてな？}{みの まわりの ふしぎ}{s1}
  \bk@howstepone{36}{0}{yattemiyou}{bkGreen}{やってみよう}{てを うごかして しらべよう}{s2}
  \bk@howstepone{72}{0}{doguu}{bkNavy}{\R{数学}{すう|がく}の \R{道具}{どう|ぐ}}{あたらしい かんがえかた}{s3}
  \bk@howstepone{108}{0}{reidai}{bkBlue}{れいだい}{ときかたの おてほん}{s4}
  \bk@howstepone{0}{-51}{jibun}{bkOrange}{じぶんでとこう}{おてほんを まねして かこう}{s5}
  \bk@howstepone{36}{-51}{tsumazuki}{bkRed}{まちがえやすい\\ところ}{まちがいの わけを しろう}{s6}
  \bk@howstepone{72}{-51}{butsuri}{bkPurple}{つかわれている\\ところ}{なにに やくだつのかな？}{s7}
  \bk@howstepone{108}{-51}{check}{bkTheme}{もんだい}{\R{章}{しょう}の まとめに チャレンジ}{s8}
  \foreach \a/\b in {s1/s2,s2/s3,s3/s4,s5/s6,s6/s7,s7/s8}
    {\draw[bk arrow, bkGray] (\a.east) -- (\b.west);}
  \draw[bk arrow, bkGray, rounded corners=2mm] (s4.south) -- ++(0,-2.5) -| (s5.north);
\end{tikzpicture}
\end{bkfigure}

\bksubsection{2}{この \R{本}{ほん}には、 かきこんで いいよ}

「じぶんでとこう」と 「もんだい」には、 かきこむ ところが あります。 えんぴつで、 しきと こたえを かこう。
\begin{bkfigure}
\begin{tikzpicture}[x=1mm, y=1mm]
  \fill[white, draw=bkOrange, line width=1pt, rounded corners=2.5mm] (0,0) rectangle (104,20);
  \node[anchor=north west, font=\sffamily\footnotesize, text=bkGray] at (3,19.5) {れい：ふうせん 3こと 2こ。 あわせて なんこ？};
  \node[anchor=west, fill=bkOrange, text=white, rounded corners=1.2mm, font=\sffamily\bfseries\small, inner xsep=4pt, inner ysep=2pt] (k) at (3,6.5) {しき};
  \node[anchor=west, font=\fontsize{15}{16}\selectfont\bfseries, text=bkBlue!70!black] at (18,6.8) {3 ＋ 2 ＝ 5};
  \node[anchor=east, font=\sffamily\bfseries\small, text=bkOrange] at (79,6.5) {こたえ};
  \draw[bkOrange, line width=1pt, rounded corners=1.5mm] (81,2) rectangle (100,11);
  \node[font=\fontsize{15}{16}\selectfont\bfseries, text=bkBlue!70!black] at (90.5,6.7) {5こ};
  \bkpencil{(110,4)}{20}{bkRoute4}
\end{tikzpicture}
\end{bkfigure}

\clearpage
\bksubsection{3}{こまったら}

\begin{enumerate}[label={\bkcircled[bkNavy]{\arabic*}}, leftmargin=2.4\zw, itemsep=4pt, topsep=2pt]
\item 「れいだい」に もどろう。 「もんだい」と おなじ かたちの もんだいを、 れいだいで といて みせて います。
\item 「ふりかえり」の 「→\,ページ」を みよう。 わからなかった ところに もどれます。
\item こたえは、 この \R{本}{ほん}の うしろの 「こたえ」の ページに あります。 \R{自分}{じ|ぶん}で といてから あわせよう。
\item \R{章}{しょう}の はじめの 「\R{道}{みち}しるべ」に、 \R{前}{まえ}に よむと よい \R{章}{しょう}が あります。
\end{enumerate}

\bksubsection{4}{この \R{本}{ほん}の 3つの やくそく}

\begin{bkfigure}
\begin{tikzpicture}[x=1mm, y=1mm, yk/.style={draw=bkTheme, line width=1.1pt, fill=white, rounded corners=3mm, minimum width=40mm, minimum height=46mm, anchor=north}]
  \foreach \i/\x in {1/0,2/46,3/92} {\node[yk] (y\i) at (\x+20,0) {};
    \node[circle, fill=bkTheme, text=white, font=\sffamily\bfseries, minimum size=7mm, inner sep=0pt] at (\x+5.5,-5.5) {\i};}
  % 1 たんい
  \foreach \x in {7,14,21} {\bkapple{(\x,-17)}{3.2}{bkAppleRed}}
  \node[font=\fontsize{15}{16}\selectfont\bfseries, text=bkInk] at (32,-17) {3こ};
  \node[anchor=north, text width=36mm, align=center, font=\fontsize{11}{17}\selectfont] at (20,-25) {かずには 「こ」「\R{本}{ほん}」 「まい」などを つけよう};
  % 2 ゆび
  \foreach \k in {0,...,4} {\bkohajiki{(53+\k*6.5,-17)}{2.6}{bkRoute6}}
  \node[anchor=north, text width=36mm, align=center, font=\fontsize{11}{17}\selectfont] at (66,-25) {ゆびや おはじきを つかって いいよ};
  % 3 まちがい
  \node[font=\fontsize{20}{20}\selectfont, text=bkRed] at (106,-16) {\bkbatsu};
  \node[font=\fontsize{12}{12}\selectfont, text=bkGray] at (112,-16) {→};
  \node[font=\fontsize{20}{20}\selectfont, text=bkBlue] at (119,-16) {\bkmaru};
  \node[anchor=north, text width=36mm, align=center, font=\fontsize{11}{17}\selectfont] at (112,-25) {まちがえても だいじょうぶ。 わけを しらべよう};
\end{tikzpicture}
\end{bkfigure}

\bksubsection{5}{3\R{人}{にん}の なかまが いるよ}

\noindent\begin{minipage}[c]{0.45\linewidth}\centering
\begin{tikzpicture}[x=1mm, y=1mm]
  \bkChar{sora}{(8,0)}{29mm}\bkChar{mina}{(25,0)}{29mm}\bkChar{delta}{(43,0)}{24mm}
\end{tikzpicture}
\end{minipage}\hfill
\begin{minipage}[c]{0.52\linewidth}
ソラと ミナと ロボットの デルタが、 ふきだしで ヒントを くれます。 こまったら、 3\R{人}{にん}の ことばを よんで みよう。
\end{minipage}
\par
\endgroup}
\makeatother
